\documentclass[onecolumn, draftclsnofoot,10pt, compsoc]{IEEEtran}
\usepackage{graphicx}
\usepackage{url}
\usepackage{setspace}

\usepackage{geometry}
\geometry{textheight=9.5in, textwidth=7in}

% 1. Fill in these details
\def \CapstoneTeamName{		The Cleverly Named Team}
\def \CapstoneTeamNumber{		45}
\def \GroupMemberOne{			Kaz Buckley}
\def \GroupMemberTwo{			Christopher Hebert}
\def \GroupMemberThree{			Aidan O'Malley}
\def \CapstoneProjectName{		Interactive Music Theory Application}
\def \CapstoneSponsorCompany{	Western Oregon University/ NA}
\def \CapstoneSponsorPerson{		Lukas Hein}

% 2. Uncomment the appropriate line below so that the document type works
\def \DocType{		Problem Statement
				%Requirements Document
				%Technology Review
				%Design Document
				%Progress Report
				}
			
\newcommand{\NameSigPair}[1]{\par
\makebox[2.75in][r]{#1} \hfil 	\makebox[3.25in]{\makebox[2.25in]{\hrulefill} \hfill		\makebox[.75in]{\hrulefill}}
\par\vspace{-12pt} \textit{\tiny\noindent
\makebox[2.75in]{} \hfil		\makebox[3.25in]{\makebox[2.25in][r]{Signature} \hfill	\makebox[.75in][r]{Date}}}}
% 3. If the document is not to be signed, uncomment the RENEWcommand below
%\renewcommand{\NameSigPair}[1]{#1}

%%%%%%%%%%%%%%%%%%%%%%%%%%%%%%%%%%%%%%%
\begin{document}
\begin{titlepage}
    \pagenumbering{gobble}
    \begin{singlespace}
    	\includegraphics[height=4cm]{coe_v_spot1}
        \hfill 
        % 4. If you have a logo, use this includegraphics command to put it on the coversheet.
        %\includegraphics[height=4cm]{CompanyLogo}   
        \par\vspace{.2in}
        \centering
        \scshape{
            \huge CS Capstone \DocType \par
            {\large\today}\par
            \vspace{.5in}
            \textbf{\Huge\CapstoneProjectName}\par
            \vfill
            {\large Prepared for}\par
            \Huge \CapstoneSponsorCompany\par
            \vspace{5pt}
            {\Large\NameSigPair{\CapstoneSponsorPerson}\par}
            {\large Prepared by }\par
            Group\CapstoneTeamNumber\par
            % 5. comment out the line below this one if you do not wish to name your team
            \CapstoneTeamName\par 
            \vspace{5pt}
            {\Large
                \NameSigPair{\GroupMemberOne}\par
                \NameSigPair{\GroupMemberTwo}\par
                \NameSigPair{\GroupMemberThree}\par
            }
            \vspace{20pt}
        }
        \begin{abstract}
        % 6. Fill in your abstract    
        	This document describes the objectives, features, metrics, and other data related to the what and why we are creating this Music Theory Application. The objective of this project is to assist musicians and composers in understanding the basics of music theory, in order to help them in their journey of creating music. Our client, Lukas, came up with the idea to create this app because of the necessity for such a tool when teaching his musical students or when attempting to understand how to create “objectively good” music. Since “good music“ is a very subjective term I define “objectively good” music as music that is aesthetically pleasant to the majority of the populous. I will go more in depth to what this means later. 
        \end{abstract}     
    \end{singlespace}
\end{titlepage}
\newpage
\pagenumbering{arabic}
\tableofcontents
% 7. uncomment this (if applicable). Consider adding a page break.
%\listoffigures
%\listoftables
\clearpage

% 8. now you write!
\section{Definition / Description}
	The current problem as described by Lukas is that students and composers often don't understand why what they play/ write sounds good.  Musicians often follow instructions of what to play without questioning why they are playing the notes in that order or why playing the notes in a specific order might sound better than playing the same notes in a different order. Likewise composers might have an understanding of how to compose within a specific key but not understand how to move chord progressions through that key in a logical and aesthetic arrangement. In addition Lukas feels that the current system for teaching this Music Theory is flawed because it tends to miss giving significance to concise and easily understandable patterns to music theory. This makes it difficult for composers and musicians to understand why they make music in a specific manner. The goal of this project is to create an interactive module that can teach users to understand some basic musical theory while also serving as a tool for users who wish to use these details to create original music.

\section{Solution / Features}
	The proposed solution is an interactive mobile application that will display detailed relationships between notes based on the key of your choice. From here you will be able to view information on how to use that key to create “objectively good” music. Everyone's opinion of what music is good, is very different, however I would argue that among a group of people with different tastes in music that they would be able to objectively judge, as a group (meaning that the majority would agree), if a specific song is aesthetic. Meaning that for any given song the group can decide if it is or isn't pleasant to listen to or could be considered music, regardless of any single persons taste in music. With this definition of a “objectively good song” one of the goals of this application would be to show that there are logical patterns that can determine if a song would meet this criteria without needing a jury of listeners. By creating and using rules for different genres of music, we can determine if any collection of random notes could be considered a song.

\section {Metrics}
	The first and most important metric for measuring the completeness of this project will be comparing our progress of the interactive application with the a list of desired features for the application. These include; the interactive circle of fifths, the logical progression of chords, and the general rules for creating western popular-culture music. These rules are spesific to the music found in often found kAmerican and western culture music This specific genre will be the focus of this application although it may be a stretch goal to incorporate additional rules for other genres of music. Additional metrics might come in the form of measuring the users knowledge of music theory through interactive media quizzes. 

\end{document}